\section{Related Work}
\label{sec:related}

\paragraph{Lightweight vision--language models.}
A VLM combines an image encoder with an LLM so as to understand and generate
visual and linguistic information jointly. Early multimodal learning centred on
contrastive learning over image--text pairs to learn a shared embedding space; as
the instruction-following capability of LLMs advanced, the dominant approach
shifted to projecting image features into the input token space of an LLM and
using instruction tuning to acquire visual question answering and description
generation abilities. LLaVA \citep{liu2023llava} is a representative case in
which an open-source LLM was fine-tuned on visual instruction data generated by
GPT-4, and many subsequent VLM studies inherited this instruction-tuning
paradigm. Commercial large-scale VLMs such as GPT-4V and Gemini secure
general-purpose visual understanding through vast parameter counts and training
data, but their inference cost and deployment constraints limit their use on
consumer-grade hardware. The four models targeted here are all lightweight VLMs
in the 2--3B class released in 2025--2026, each pursuing parameter efficiency in
a different way. Qwen2.5-VL \citep{bai2025qwen25vl} is characterized by
dynamic-resolution processing, which adjusts the number of visual tokens
according to image size, and by an MRoPE extension that aligns temporal
information to absolute time. Its successor, Qwen3-VL, expanded into dense
(2B/4B/8B/32B) and mixture-of-experts (30B-A3B/235B-A22B) families and introduced
a ``thinking mode'' together with DeepStack-style multi-level fusion of visual
features. SmolVLM2 \citep{marafioti2025smolvlm} is an on-device-oriented
architecture combining a SigLIP image encoder with the SmolLM2 language model,
focusing on extremely low memory usage (5.2\,GB for video inference in the 2.2B
model). Gemma4-E2B is a mixture-of-experts-family structure that activates only
2.3B parameters at inference through per-layer embeddings (PLE) while exploiting
the expressive capacity of a 5.1B-class total parameter count.

\paragraph{Parameter-efficient fine-tuning.}
LoRA \citep{hu2022lora} freezes a pretrained weight matrix $W_0$ and learns an
approximation of its update $\Delta W$ as the product $BA$ of two low-rank
matrices $A \in \mathbb{R}^{r \times k}$ and $B \in \mathbb{R}^{d \times r}$, with
$r \ll \min(d,k)$. The forward pass is $h = W_0 x + BAx$, and since only $A$ and
$B$ are trained, the ratio of trainable to total parameters can be reduced
dramatically. The rank $r$ governs the trade-off between expressive capacity and
parameter efficiency, while alpha scales $\Delta W$ (multiplying $BAx$ by
$\alpha/r$). QLoRA \citep{dettmers2023qlora} combines LoRA with 4-bit
quantization to reduce VRAM usage further, through (1) 4-bit NormalFloat (NF4)
quantization optimized for normally distributed weights, (2) double quantization,
which quantizes the quantization constants themselves, and (3) a paged optimizer
that offloads GPU memory spikes to the CPU; the original paper showed that this
allows fine-tuning of a 65B-class model on a single 48\,GB GPU while maintaining
performance close to full 16-bit fine-tuning. Other PEFT methods include partial
fine-tuning, prefix tuning, and adapter insertion between transformer layers.
Compared with these, the advantages of the LoRA family are that $BA$ can be
merged into $W_0$ at inference time so that no additional latency is incurred,
and that adapters can be stored and swapped independently, allowing a single base
model to be reused across multiple tasks. We adopt LoRA/QLoRA because these two
advantages are particularly favourable for the repeated hyperparameter search
conducted on consumer GPUs.

\paragraph{Medical VQA benchmarks and medical-specialised models.}
Medical VQA is the task of generating a correct answer given a medical image (a
pathology tissue slide, a radiological image, and so on) together with a
natural-language question; unlike general-domain VQA it requires both an
understanding of specialist medical terminology and fine-grained visual grounding
of imaging findings. Question types divide broadly into closed-ended questions
answered yes/no or by selecting among options and open-ended questions requiring
free-form answers, and the two types differ greatly in scoring method and
difficulty. PathVQA \citep{he2020pathvqa} provides 32,799 question--answer pairs
over 4,998 pathology tissue images extracted from pathology textbooks and the
PEIR digital library, organized into seven question types modelled on the
American Board of Pathology certification examination format. SLAKE
\citep{liu2021slake} combines 14,028 bilingual English--Chinese question--answer
pairs over 642 radiology/CT images with a medical knowledge graph of 5,232
knowledge triplets. VQA-RAD \citep{lau2018vqarad} was the first dataset to
collect questions that clinicians naturally raised while viewing radiological
images (CT/MRI/X-ray), containing roughly 3,500 question--answer pairs over 315
images. Attempts to specialize general-purpose VLMs for the medical domain divide
into large-scale retraining and adaptation approaches. LLaVA-Med
\citep{li2023llavamed} generated GPT-4-based instruction data from large-scale
biomedical figure--caption data in PubMed Central and adapted general-purpose
LLaVA to the biomedical domain through curriculum learning. Med-Flamingo
\citep{moor2023medflamingo} performed continued pretraining of OpenFlamingo-9B on
image--text data from medical papers and textbooks, acquiring the ability to
adapt to medical VQA from only a few examples. CheXagent
\citep{chen2024chexagent} is a foundation model composed of a clinical LLM, a
visual encoder, and a cross-modal bridge network specialized for chest X-ray
interpretation. These prior studies commonly target models on the scale of
billions to hundreds of billions of parameters and presuppose either continued
pretraining on large biomedical corpora or the generation of large-scale
instruction data; by contrast, we fine-tune 2--3B-class lightweight models with
QLoRA using only small domain datasets (thousands to tens of thousands of
samples).

\paragraph{Hyperparameter optimization, from random search to LLM agents.}
The simplest HPO method is grid search, whose required number of evaluations
grows exponentially with the dimensionality of the search space. Bergstra and
Bengio \citep{bergstra2012random} showed theoretically and empirically that
random sampling of combinations can yield substantially better --- or equivalent
--- performance than grid search within the same computational budget, because
the effective dimensionality of the hyperparameters that actually influence
performance is often low. Bayesian optimization constructs a probabilistic
surrogate model of the objective from previous evaluations and selects the next
evaluation point accordingly. Optuna \citep{akiba2019optuna}, used as the control
here, adopts the tree-structured Parzen estimator (TPE) algorithm and provides a
define-by-run API together with efficient pruning strategies. Among
early-stopping-based methods, Hyperband \citep{li2018hyperband} formulates HPO as
a pure-exploration bandit problem by repeatedly applying successive halving,
allocating little resource to unpromising configurations and progressively more
to promising ones; for this to hold there must be a significant rank correlation
between the early learning curve and final converged performance. In contrast to
Bayesian optimization, which relies on a numerical surrogate model, a scheme in
which an LLM agent directly interprets previous experimental logs and proposes
the next configuration through natural-language reasoning has recently been
discussed as a new alternative. Its theoretical distinctions are three: the LLM
can transfer knowledge across different tasks and architectures by drawing on
domain knowledge acquired during pretraining; unlike Bayesian optimization, which
treats individual hyperparameters as independent variables, it can understand
interactions among hyperparameters structurally on the basis of prior knowledge;
and it records the reason for changing a configuration explicitly in natural
language, providing interpretability and traceability. This is, however, less an
established standard methodology than one instance of an emerging research
current on the use of LLM agents for automating scientific experimentation.

\paragraph{What distinguishes this work.}
Research on medical-specialized VLMs generally presupposes large models and large
domain datasets; research on PEFT methodology concentrates on verifying parameter
efficiency in general domains; and research on HPO automation has rarely compared
LLM-agent-based methods directly against established Bayesian optimization on the
same task. At the intersection of these three currents --- consumer GPU
environment, lightweight VLMs, QLoRA domain adaptation, and LLM-based autonomous
HPO --- this paper verifies the zero-shot and fine-tuned medical VQA performance
of lightweight VLMs with statistical rigour (RQ1, RQ2) under a single protocol
spanning all three stages, and validates the practical value of
autoresearch-style autonomous search by comparing it directly against Optuna
(TPE) (RQ3), reporting a negative result together with a verified failure
mechanism.
